Due to the precession of the Earth's axis, the region of sky visible from a
location with fixed geographical coordinates changes with time. Is it possible
that, at some point in time, Sirius will not rise as seen from Krakow, while
Canopus will rise and set? Assume that the Earth's axis traces out a cone of
angle $47^\circ$. Krakow is at latitude $50.1^\circ$\,N; the current
equatorial coordinates (right ascension and declination) of these stars are:
\[
\begin{array}{lll}
\text{Sirius } (\alpha\,\text{CMa}): &
  6^{\mathrm{h}}45^{\mathrm{m}}, & -16^\circ 43' \\
\text{Canopus } (\alpha\,\text{Car}): &
  6^{\mathrm{h}}24^{\mathrm{m}}, & -52^\circ 42'
\end{array}
\]
